\section{Discussion}
\label{sec:discussion}

\subsection{What this paper accomplishes}

\paragraph{Breaks the Wamura problem.}
Theorem~\ref{thm:convergence} shows that controlled relaxation
drives risk estimates to truth at exponential rate, converting
``permanent overcaution'' into ``temporary overcaution with a
characterized resolution time.''  This closes the most important
named open problem of~\citep{lasser2026horizon}.

\paragraph{Provides a damage-bounded mechanism.}
Theorem~\ref{thm:damage_bound} bounds the maximum damage from a
test by the user-controllable parameter $\damagebound$.
Controlled relaxation is therefore safe to deploy: the cost of
learning is bounded even when the risk turns out to be real, and
the cooperative-closure condition of
Definition~\ref{def:test_scope}(ii) prevents cross-boundary
cascades through shared cooperative capabilities.

\paragraph{Integrates with the existing machinery.}
Propositions~\ref{prop:preserve_balance},
\ref{prop:empirical_arbitration} and Lemma~\ref{lem:sixth_channel}
show that the protocol preserves the self-balancing property of
$\volL$, resolves inter-agent risk disagreements, and extends
the multi-channel detection architecture.  The protocol is
additive to the existing framework, not a replacement.

\paragraph{Resolves the Loop 3 phase-boundary degeneracy.}
The degenerate phase boundary $\rS = 0$ for Loop~3
in~\citep[Proposition~4]{lasser2026phase} arises precisely
because the Wamura asymmetry eliminates the exogenous
correction rate on the restricted dimension.  Controlled
relaxation makes $\rS > 0$ for Loop~3 (via the exercise-regime
observations of Section~\ref{sec:relaxation_protocol}),
restoring the phase boundary to the non-degenerate form that
the design criterion
of~\citep[Theorem~2]{lasser2026phase} requires for
self-correction.  The two papers are complementary: one
characterizes when self-correction works in aggregate; this one
supplies the micro-mechanism that makes self-correction work
for the restriction-specific feedback loop.

\subsection{Limitations}

\paragraph{B-to-C gap inheritance.}
All results are stated in terms of $\volL$.
Section~\ref{sec:framing_precondition} records that $\damagebound$
bounds contraction of $\volL$, not of the experiential target
directly; if the gap between them is large, even a
controlled-relaxation-clean restriction landscape need not
correspond to experiential optimality.  The gap-closure programme
is a separate research direction, noted
in~\citep{lasser2026gfm} and outside the scope of this paper.

\paragraph{Honest-channel assumption.}
We assume honest observation channels throughout.  Adversarial
manipulation of the detection signal---a compromised channel
reporting fabricated contraction or concealing genuine
contraction---falls outside the scope of the proof and is the
domain of the commitment protocols of~\citep{lasser2026exo}.
Controlled relaxation composes cleanly with those
protocols---the ledger of~\citep[Section~5]{lasser2026exo} can
record test outcomes as committed observations---but the
composition itself is not formalized here.

\paragraph{Identifiability under misspecification.}
Assumption~\ref{ass:T2} requires the Bayesian observation model
to match the true data-generating process.  If $\alphafp$ or
$\beta$ are misspecified, Theorem~\ref{thm:convergence}'s
log-likelihood-ratio argument still converges, but it converges
to the wrong hypothesis: under the misspecified model, the
posterior concentrates on whichever hypothesis is closer in
KL-divergence to the true distribution, not on the true
hypothesis.  Detector calibration is therefore a prerequisite
for the convergence guarantee to carry its safety
interpretation.

\subsection{Open questions}

\paragraph{Damage-tolerance calibration.}
The tolerance $\damagebound$ is a free parameter that trades off
learning speed against risk exposure.  Options include
\begin{itemize}
\item a fixed fraction of $\volL$ (e.g.,
  $\damagebound = 0.01 \cdot \volL$);
\item proportional to restriction cost
  ($\damagebound = c \cdot |\Delta \volL(\Xirestrict)|$);
\item proportional to information value
  ($\damagebound = c' \cdot \IV(T)$);
\item user-set per test (maximum safety, maximum user burden).
\end{itemize}
The worked example (Section~\ref{sec:worked_example}) shows that
even small $\damagebound$ produces rapid convergence when
detection sensitivity is high, so the answer may depend more on
the detection architecture than on $\damagebound$ itself.

\paragraph{Transient versus persistent pathways.}
Section~\ref{sec:detection_sensitivity} assumes persistent
pathway activation (Assumption~\ref{ass:D0}).  Some risk
pathways may be transient---they produce a brief contraction
that the detector could miss at low sampling rate.  The current
formulation handles transient pathways via
$\beta(\testdur) = (1 - \deltaagg)^{\min(k, \testdur)}$ where
$k$ is the pathway duration, but if $k = 1$ then
$\beta = 1 - \deltaagg$ may be non-negligible.  Whether the
protocol should require longer test windows for suspected
transient pathways, or whether the detection architecture should
raise its sampling rate during tests, is an engineering choice
left open here.

\paragraph{Adversarial risk claimants.}
Assumption~\ref{ass:T3}(a) rules out adaptation by the tested
agent.  What happens when an agent strategically inflates risk
claims to lock competitors' capabilities?  The controlled
relaxation protocol still works (it generates evidence
regardless of the claimant's intent), but the meta-policy for
test initiation may need to be more aggressive when strategic
inflation is suspected.  Detection of strategic inflation is a
variant of the scorpion-detection
problem~\citep[Propositions~3--4]{lasser2026scm}: an agent
whose risk claims are systematically disconfirmed by
controlled-relaxation tests should see its risk-trust
$\Trisk_j$ decrease, eventually down-weighting its future
claims in aggregation.  Whether the meta-policy should
explicitly model this adversarial case, or whether the
trust-dynamics response is sufficient, is an open design
question.

\subsection{Relationship to prior and subsequent papers}

\paragraph{Prior papers.}
The protocol sits downstream of and draws on:
\begin{itemize}
\item \citet{lasser2026poset} for axiomatic properties of
  $\volL$ preserved under relaxation;
\item \citet{lasser2026horizon} for the risk machinery (exercise
  pathways, concentration risk, restriction criterion) the
  protocol operates within;
\item \citet{lasser2026scm} for the multi-channel detection
  architecture that supplies $\deltaagg$ and the log-odds
  aggregation into which test outcomes feed;
\item \citet{lasser2026exo} for the verification-asymmetry
  framing and the cross-substrate isolation that justifies the
  conservative detection bound
  (Proposition~\ref{prop:detection_bounds}(c));
\item \citet{lasser2026phase} for the dynamical-systems context:
  Loop~3's degenerate phase boundary is the problem this paper
  solves.
\end{itemize}

\paragraph{Forward connections.}
Controlled relaxation generates exercise events that increase
realized-capability volume $\vol_{\mathrm{R}}$: when a capability
is locked by a false risk claim, $\vol_{\mathrm{R}}$ is
artificially suppressed, and lifting the lock lets the
capability contribute to $\vol_{\mathrm{R}}$.  The convergence
of the restriction landscape
(Proposition~\ref{prop:preserve_balance}(c)) implies that
$\vol_{\mathrm{R}} / \volL$ approaches its true value as false
restrictions clear, partially closing the B-to-C gap for the
class of capabilities locked by the Wamura problem.  A full
treatment of this partial closure is beyond the scope of this
paper.

The protocol also provides a clean experimental signal for
structure learning: the SCM
of~\citep[Definition~1]{lasser2026scm} can use designed-experiment
observations from controlled relaxation (where the intervention
is known and the outcome is observed) as ground-truth data
points for calibrating structural confidence.  This is strictly
more informative than passive observational data and could
accelerate SCM calibration beyond the rates achievable under
passive observation alone.
